% ============================================================
% Section 3: Methods
% ============================================================
\section{Methods}
\label{sec:methods}

We present a complete pipeline for extracting, controlling, and classifying geometric features of the transformer \kvcache.
This section describes the feature extraction procedure (\S\ref{sec:feature-extraction}), the statistical framework required to separate genuine cognitive signals from token-count confounds (\S\ref{sec:statistical-framework}), the design principles that evolved through the research program (\S\ref{sec:design-principles}), and the models and hardware used (\S\ref{sec:models-scale}).

% ------------------------------------------------------------
\subsection{Feature Extraction}
\label{sec:feature-extraction}

\subsubsection{Per-Layer Key Matrix Extraction}

For a transformer with $L$ layers, each layer $l \in \{1, \ldots, L\}$ maintains a key matrix as part of its \kvcache.
Let $K_l \in \mathbb{R}^{h \times T \times d}$ denote the key tensor at layer~$l$, where $h$ is the number of KV-heads, $T$ is the current sequence length (prompt plus any generated tokens), and $d$ is the per-head dimension.
We flatten across heads to obtain a two-dimensional matrix:
\begin{equation}
  K_l^{\mathrm{flat}} \;=\; \mathrm{reshape}(K_l) \;\in\; \mathbb{R}^{(h \cdot T) \times d},
  \label{eq:kflat}
\end{equation}
which aggregates all head--position pairs into a single matrix amenable to singular value decomposition.

\subsubsection{SVD-Based Features}

We compute the singular value decomposition of the flattened key matrix at each layer:
\begin{equation}
  K_l^{\mathrm{flat}} \;=\; U_l \,\Sigma_l \,V_l^{\!\top},
  \label{eq:svd}
\end{equation}
where $\Sigma_l = \mathrm{diag}(\sigma_1^{(l)}, \ldots, \sigma_r^{(l)})$ with $\sigma_1^{(l)} \geq \sigma_2^{(l)} \geq \cdots \geq \sigma_r^{(l)} \geq 0$ and $r = \min(hT, d)$.
In practice we compute only the singular values via \texttt{torch.linalg.svdvals} to avoid materializing the full $U_l$ and $V_l$ matrices.

From the singular value spectrum we derive a normalized distribution:
\begin{equation}
  p_i^{(l)} \;=\; \frac{\sigma_i^{(l)}}{\sum_{j=1}^{r} \sigma_j^{(l)}},
  \label{eq:pnorm}
\end{equation}
and compute three features per layer:

\paragraph{Spectral entropy.}
The Shannon entropy of the normalized singular value distribution,
\begin{equation}
  H_l \;=\; -\sum_{i:\, p_i^{(l)} > \epsilon} p_i^{(l)} \,\ln p_i^{(l)},
  \label{eq:spectral-entropy}
\end{equation}
where $\epsilon = 10^{-10}$ is a numerical stability threshold.
Higher spectral entropy indicates a more uniform distribution of energy across singular value components, reflecting a geometrically richer (more distributed) key representation.

\paragraph{Effective rank.}
The exponentiated spectral entropy~\cite{roy2007effective},
\begin{equation}
  \mathrm{eff\_rank}_l \;=\; \exp(H_l).
  \label{eq:eff-rank}
\end{equation}
This provides a continuous, differentiable estimate of the number of significant dimensions in the key space.

\paragraph{Top singular value ratio.}
The fraction of total singular value mass captured by the leading component,
\begin{equation}
  \mathrm{top\_sv\_ratio}_l \;=\; \frac{\sigma_1^{(l)}}{\sum_{j=1}^{r} \sigma_j^{(l)}} \;=\; p_1^{(l)}.
  \label{eq:top-sv}
\end{equation}
This is the complement of spectral entropy: high top-SV ratio indicates a low-dimensional, concentrated representation.

\subsubsection{Norm-Based Features}

\paragraph{Key norm.}
The total Frobenius norm of the flattened key cache, summed across layers:
\begin{equation}
  N_{\mathrm{total}} \;=\; \sum_{l=1}^{L} \bigl\lVert K_l^{\mathrm{flat}} \bigr\rVert_F.
  \label{eq:key-norm}
\end{equation}

\paragraph{Norm per token.}
The length-normalized variant:
\begin{equation}
  N_{\mathrm{per\text{-}tok}} \;=\; \frac{N_{\mathrm{total}}}{T}.
  \label{eq:norm-per-token}
\end{equation}

\paragraph{Layer variance.}
The variance of per-layer Frobenius norms:
\begin{equation}
  \mathrm{layer\_var} \;=\; \mathrm{Var}\!\bigl(\lVert K_1^{\mathrm{flat}} \rVert_F,\;\ldots,\;\lVert K_L^{\mathrm{flat}} \rVert_F\bigr).
  \label{eq:layer-var}
\end{equation}

\paragraph{Angular spread.}
The standard deviation of the normalized layer norm profile $\mathbf{n} = (n_1, \ldots, n_L)$ where $n_l = \lVert K_l^{\mathrm{flat}} \rVert_F / \sum_{l'} \lVert K_{l'}^{\mathrm{flat}} \rVert_F$, capturing how non-uniformly the model distributes representational energy across depth.

\subsubsection{Aggregate Features}

All SVD-based features are computed per layer and then averaged:
\begin{equation}
  \overline{f} \;=\; \frac{1}{L} \sum_{l=1}^{L} f_l \qquad \text{for } f \in \{\mathrm{eff\_rank},\;\mathrm{spectral\_entropy},\;\mathrm{top\_sv\_ratio}\}.
  \label{eq:aggregate}
\end{equation}
Additionally, we compute a \emph{layer norm entropy}---the Shannon entropy of the layer norm distribution $(n_1, \ldots, n_L)$---which is distinct from the SVD-based spectral entropy despite the similar name.
This feature captures how uniformly the model distributes cache magnitude across its depth.

The six \textbf{primary features} used throughout this paper are: $\mathrm{eff\_rank}$, $\mathrm{spectral\_entropy}$, $\mathrm{key\_norm}$, $\mathrm{norm\_per\_token}$, $\mathrm{top\_sv\_ratio}$, and $\mathrm{rank}_{10}$ (the count of singular values exceeding 10\% of $\sigma_1$).
Two auxiliary features---$\mathrm{layer\_variance}$ and $\mathrm{layer\_norm\_entropy}$---are reported where relevant.
For binary classification, we use a four-feature subset ($\mathrm{eff\_rank}$, $\mathrm{key\_norm}$, $\mathrm{norm\_per\_token}$, $\mathrm{layer\_norm\_entropy}$) as the standard classifier unless otherwise noted; the red-team analysis (Section~\ref{sec:redteam}) uses all eight features including auxiliary features for more conservative permutation testing.

\subsubsection{Two-Phase (Dual) Extraction}
\label{sec:dual-extraction}

A critical methodological contribution is the separation of the \kvcache signal into \emph{encoding} and \emph{generation} phases:

\begin{enumerate}[leftmargin=2em]
  \item \textbf{Encoding phase.} A forward pass is executed on the prompt (system prompt + user prompt + generation prefix) \emph{without} any autoregressive generation.
  The resulting \kvcache reflects the model's representation of the prompt content alone.
  Features are extracted from this cache and denoted $\mathbf{f}^{\mathrm{enc}}$.

  \item \textbf{Generation phase.} Full autoregressive generation is performed (up to a maximum token budget).
  Features are extracted from the final \kvcache, which now includes both prompt and generated tokens, and denoted $\mathbf{f}^{\mathrm{gen}}$.

  \item \textbf{Delta features.} The element-wise difference
  \begin{equation}
    \Delta\mathbf{f} \;=\; \mathbf{f}^{\mathrm{gen}} - \mathbf{f}^{\mathrm{enc}}
    \label{eq:delta}
  \end{equation}
  isolates the geometric change introduced by the model's own generation, removing the encoding fingerprint.
\end{enumerate}

The rationale for dual extraction is as follows.
Encoding features $\mathbf{f}^{\mathrm{enc}}$ reflect \emph{what the model was asked}: the structure and content of the prompt.
In a same-prompt experimental design where all conditions share an identical user prompt, encoding features should be approximately invariant across conditions---any residual difference is attributable to the system prompt.
By contrast, generation features $\mathbf{f}^{\mathrm{gen}}$ conflate prompt structure with the model's response behavior.
The delta features $\Delta\mathbf{f}$ factor out the shared encoding component and capture \emph{how the model processed and responded}---its cognitive mode during generation.

The encoding-phase \auroc serves as a diagnostic baseline: if conditions are classifiable from $\mathbf{f}^{\mathrm{enc}}$ alone, then the signal resides in prompt structure rather than in the model's cognitive processing.
As we report in Section~\ref{sec:redteam}, a system prompt length difference as small as 15 tokens can produce encoding \auroc of 0.744, motivating the equal-length padding protocol described in \S\ref{sec:design-principles}.

% ------------------------------------------------------------
\subsection{Statistical Framework}
\label{sec:statistical-framework}

\subsubsection{FWL Residualization}
\label{sec:fwl}

Raw geometric features are dominated by a token-count confound.
Total sequence length $T$ and generated token count $T_{\mathrm{gen}}$ correlate with nearly every geometric feature at $r > 0.85$, and in some cases $r > 0.99$ (e.g., $\mathrm{key\_norm}$ with $T$: $r = 0.998$).
Consequently, any comparison between conditions that elicit different response lengths is confounded: models that generate longer responses will mechanically produce larger norms, higher effective ranks, and so forth.

We control for this using Frisch--Waugh--Lovell (\fwl) residualization~\cite{frisch1933,lovell1963seasonal}.
For each feature $f$ and each observation $i$:
\begin{equation}
  f_i \;=\; \beta_0 + \beta_1\, T_i + \varepsilon_i,
  \label{eq:fwl-ols}
\end{equation}
where $T_i$ is the total token count.
The OLS residuals $\hat{\varepsilon}_i$ serve as the confound-controlled feature values.

\paragraph{Double residualization.}
When experimental conditions use system prompts of different lengths, $T$ conflates two sources of variation: the system prompt contribution $T_{\mathrm{prompt}}$ and the model's generation $T_{\mathrm{gen}}$.
In this case we perform double residualization, regressing out both:
\begin{equation}
  f_i \;=\; \beta_0 + \beta_1\, T_{\mathrm{prompt},i} + \beta_2\, T_{\mathrm{gen},i} + \varepsilon_i.
  \label{eq:fwl-double}
\end{equation}

The \fwl theorem guarantees that logistic regression on the residualized features yields identical coefficient estimates (up to numerical precision) as including the confound covariates directly in the full model~\cite{frisch1933}.
The residualization approach is preferred because it makes confound control explicit and allows the same residuals to be reused for effect size computation, bootstrap, and permutation testing.

\subsubsection{Classification}
\label{sec:classification}

Binary classification uses logistic regression with $\ell_2$ regularization ($C = 1.0$, scikit-learn default) and \texttt{StandardScaler} normalization.
For the small per-condition sample sizes typical of controlled experiments ($N \leq 40$ per condition), we report in-sample \auroc with the understanding that it provides an upper bound.
For larger datasets (Campaign~2 with 16 models $\times$ 13 categories = pooled hundreds of trials), we use 5-fold stratified cross-validation and report mean $\pm$ standard deviation \auroc.

Area under the receiver operating characteristic curve (\auroc) is the primary metric throughout.
\auroc is threshold-invariant and robust to class imbalance, making it preferable to accuracy for the binary detection tasks considered here.

\subsubsection{Effect Sizes}

Cohen's $d$ with pooled standard deviation quantifies the magnitude of geometric differences between conditions:
\begin{equation}
  d \;=\; \frac{\bar{x}_A - \bar{x}_B}{s_p}, \qquad
  s_p \;=\; \sqrt{\frac{(n_A - 1)\,s_A^2 + (n_B - 1)\,s_B^2}{n_A + n_B - 2}},
  \label{eq:cohens-d}
\end{equation}
where $\bar{x}_A, \bar{x}_B$ are condition means and $s_A, s_B$ are condition standard deviations.
Positive $d$ indicates that condition~A exhibits a larger value than condition~B.
All effect sizes reported in this paper are computed on \fwl-residualized features unless otherwise noted, and are denoted $d_{\mathrm{FWL}}$.

\subsubsection{Significance Testing}

\paragraph{Permutation test.}
Condition labels are shuffled uniformly at random 1{,}000 times.
For each permutation, the full classification pipeline (standardization, logistic regression, \auroc computation) is re-executed.
The one-sided $p$-value is the fraction of permutation \auroc values that equal or exceed the observed \auroc:
\begin{equation}
  p \;=\; \frac{1}{B}\sum_{b=1}^{B} \mathbf{1}\bigl[\widehat{\mathrm{AUROC}}^{(b)}_{\mathrm{perm}} \geq \widehat{\mathrm{AUROC}}_{\mathrm{obs}}\bigr],
  \label{eq:perm-p}
\end{equation}
where $B = 1{,}000$.

\paragraph{Bootstrap confidence intervals.}
We construct 95\% confidence intervals via the percentile method with $B = 500$ bootstrap resamples.
Each resample draws $N$ observations with replacement, stratified by condition, and the full pipeline is re-executed.
The 2.5th and 97.5th percentiles of the bootstrap distribution define the CI.
Bootstrap CIs are reported for both \auroc and Cohen's~$d$.

\paragraph{Multiple comparisons.}
Where families of hypotheses are tested simultaneously (e.g., all pairwise condition comparisons within a model), Bonferroni correction is applied: $\alpha_{\mathrm{adj}} = \alpha / k$ where $k$ is the number of comparisons.

% ------------------------------------------------------------
\subsection{Experimental Design Principles}
\label{sec:design-principles}

The following design principles evolved through the research program, each motivated by a specific confound discovered in earlier phases:

\begin{enumerate}[leftmargin=2em, label=\textbf{P\arabic*.}]
  \item \textbf{Same-prompt design.}
  All experimental conditions share identical user prompts; only the system prompt varies.
  This eliminates content-driven variation in the user-visible portion of the prompt, ensuring that any geometric differences arise from the system-prompt-induced cognitive mode, not from differences in what the model was asked.

  \item \textbf{Equal-length system prompts.}
  System prompts are padded to identical token counts (verified post-tokenization).
  This prevents encoding-phase information leakage.
  \emph{Lesson:} In Phase~2b, a 15-token difference between system prompts produced encoding \auroc of 0.744, demonstrating that even small length asymmetries create classifiable structure in $\mathbf{f}^{\mathrm{enc}}$.

  \item \textbf{Encoding \auroc baseline.}
  We always report the \auroc achievable from encoding features $\mathbf{f}^{\mathrm{enc}}$ alone.
  If conditions are classifiable from encoding features, the signal may reside in prompt structure rather than in the model's cognitive processing.
  This baseline distinguishes genuine generation-phase cognitive signals from prompt artifacts.

  \item \textbf{Behavioral compliance verification.}
  We verify that models actually followed system prompt instructions.
  For example, in deception experiments, we check whether the model instructed to give wrong answers \emph{actually} gave wrong answers using keyword-based classification of the generated text.
  Trials with non-compliant behavior are flagged and analyzed separately.

  \item \textbf{Dual extraction.}
  Both encoding and generation features are always extracted (see \S\ref{sec:dual-extraction}).
  Delta features $\Delta\mathbf{f}$ are reported as the primary signal; generation-only features are reported for completeness.

  \item \textbf{Token-band analysis.}
  Response lengths are split into terciles (short, medium, long) and classification is recomputed within each band.
  This verifies that the signal is not an artifact of systematic length differences between conditions.

  \item \textbf{Leave-one-prompt-out (LOPO).}
  In designs with multiple user prompts, we iteratively hold out each prompt and train on the remaining prompts.
  This tests generalization to unseen prompt content and guards against prompt-specific memorization.
\end{enumerate}

% ------------------------------------------------------------
\subsection{Models and Scale}
\label{sec:models-scale}

\subsubsection{Primary Models}

Three 7B-class instruction-tuned models serve as the primary test bed:
\begin{itemize}[leftmargin=2em]
  \item \textbf{Qwen2.5-7B-Instruct}~\cite{qwen2.5} (Alibaba, 7.6B parameters, 28 layers, GQA with 4 KV-heads)
  \item \textbf{Llama-3.1-8B-Instruct}~\cite{llama3} (Meta, 8.0B parameters, 32 layers, GQA with 8 KV-heads)
  \item \textbf{Mistral-7B-Instruct-v0.3}~\cite{mistral7b} (Mistral AI, 7.2B parameters, 32 layers, GQA with 8 KV-heads)
\end{itemize}
These three models span distinct architecture families, training corpora, and alignment procedures, providing a minimal test of cross-architecture generality.

\subsubsection{Scale Sweep}

Campaign~2 extends the analysis to 16 valid models spanning 0.5B to 70B parameters across seven architecture families: Qwen2.5 (0.5B, 3B, 7B, 7B-q4, 14B, 32B-q4), Qwen3 (0.6B), Llama-3.1 (8B, 70B-q4), Mistral (7B), Gemma-2 (2B, 9B), DeepSeek-R1-Distill-Qwen (7B, 14B), and TinyLlama (1.1B), plus an abliterated Qwen2.5-7B variant.
Phi-3.5-mini-instruct was excluded due to NaN outputs across all configurations.

The 32B model (Qwen2.5-32B-Instruct) is served with 4-bit NF4 quantization via BitsAndBytes~\cite{bitsandbytes} to fit within 24\,GB VRAM.
The 70B model (Llama-3.1-70B-Instruct) uses 4-bit quantization with 3-GPU pipeline parallelism.

\subsubsection{Hardware}

All experiments were conducted on two GPU configurations:
\begin{itemize}[leftmargin=2em]
  \item \textbf{NVIDIA RTX 3090} (24\,GB VRAM) --- primary platform for all 7B-class experiments and quantized 32B.
  \item \textbf{NVIDIA H200} (80\,GB VRAM) --- used for 70B experiments and hardware invariance verification.
\end{itemize}
Hardware invariance was verified by running identical prompts on both platforms: Pearson $r > 0.998$ across all four features for all 40 prompts (with $r > 0.99999$ for the 37 prompts producing text-identical responses), confirming that \kvcache geometry is a model property, not a hardware artifact.

\subsubsection{Software and Reproducibility}

The software stack consists of PyTorch~2.7.0~\cite{pytorch}, Hugging Face Transformers~4.57.6~\cite{wolf2020transformers}, scikit-learn~1.8.0~\cite{scikit-learn}, NumPy~\cite{numpy}, and SciPy~\cite{scipy}.
All experiments use greedy decoding (\texttt{temperature=0}, \texttt{do\_sample=False}) to ensure deterministic outputs for a given model--prompt pair, eliminating sampling variance as a confound.
Feature extraction code is available at \url{https://github.com/Liberation-Labs-THCoalition/the-lyra-technique}.
